\section{Why One-Shot Editing Fails on Gemma: Norm Placement}
\label{sec:analysis}

In a pre-norm block (Llama, Qwen, Mistral) a sublayer writes $h \leftarrow h + W_o x$, so an adapter adds exactly
$\dW x$ to the residual stream. Gemma~2/3 \citep{gemma3} (and OLMo~2 \citep{olmo2}) normalize the sublayer output first:
\begin{equation}
  h \leftarrow h + g \odot \mathrm{RMSNorm}\!\left(W_o x + \dW x\right).
\end{equation}
No adapter-writable module then writes linearly into the residual stream: an edit is (i) renormalized, so it can
redirect but not amplify the sublayer's write; (ii) distorted per channel by $g$; (iii) scaled differently for every
input. One-shot sizing from local effects therefore undershoots or, when enlarged, crowds out the sublayer's normal
output.

\paragraph{Evidence.} (1) Edit sizes: on Qwen, \astra{}'s edits are $0.16$ and $0.12\lVert W\rVert$; on Gemma, its
attention edit reaches $1.95\lVert W\rVert$ under a $2\times$ cap and $2.47\lVert W\rVert$ under a $5\times$ cap, still
recovering only 15--20\% and costing plain accuracy. (2) Saturation: scaling Gemma v1's edit by $\alpha$ (chosen on
train facts) lowers held-out follow rate only to a plateau of ${\sim}35$--$40\%$ while agreement with correct guesses
collapses (\cref{fig:scale}). (3) Gradient LoRAs succeed on both models because they optimize through the norm.

\begin{figure}[t]
\centering
\fbox{\parbox{0.8\linewidth}{\centering\vspace{3em}\todo{scale-sweep figure from results/m05/diagnose\_*.json}\vspace{3em}}}
\caption{Scaling \astra{}'s Gemma v1 edit: follow rate saturates while agreement with correct guesses collapses.}
\label{fig:scale}
\end{figure}

\paragraph{Practical rule.} Prefer pre-norm targets for one-shot LoRA editing; treat post-sublayer-norm models as a
stress test. \todo{direct test: residual-stream displacement vs.\ $\alpha$ on both models}
